% !TeX root = main.tex
\newpage
\section{Khái niệm hệ thống và cơ sở lý thuyết}
\subsection{Nền tảng IoT}
ESP32 là họ vi điều khiển 32-bit do Espressif Systems phát triển, kế thừa từ ESP8266. Nền tảng này tích hợp kết nối Wi-Fi và Bluetooth/BLE, đồng thời cung cấp hiệu năng xử lý đủ tốt cho các ứng dụng IoT biên. Trong đồ án này, ESP32-S3 được dùng làm bộ điều khiển trung tâm cho các tác vụ cảm biến, xử lý cục bộ, truyền thông và tương tác với đám mây.

\subsubsection{Thông số kỹ thuật chính}
\begin{itemize}
    \item \textbf{Vi xử lý (Processor):} Kiến trúc 32-bit Xtensa/RISC-V (tùy thuộc vào model), thường hoạt động trong dải tần số 160--240 MHz.
    \item \textbf{Bộ nhớ (Memory):} Tài nguyên SRAM tích hợp trên chip với hỗ trợ bộ nhớ Flash ngoài (thông dụng từ 4--16 MB, tùy thuộc vào cấu hình của từng bo mạch).
    \item \textbf{Kết nối (Connectivity):} Hỗ trợ Wi-Fi chuẩn IEEE 802.11 b/g/n và Bluetooth Low Energy (BLE).
    \item \textbf{Giao tiếp ngoại vi (Peripheral Interfaces):} Tích hợp các chuẩn giao tiếp đa dạng như GPIO, PWM, ADC, UART, SPI và I$^2$C.
    \item \textbf{Tính năng quản lý năng lượng (Power Features):} Hỗ trợ các chế độ tiết kiệm năng lượng (low-power modes), phù hợp cho các hệ thống nhúng và thiết bị IoT chạy bằng pin.
    \item \textbf{Điều kiện hoạt động (Operating Range):} Đáp ứng các tiêu chuẩn hoạt động trong môi trường công nghiệp cho các kịch bản triển khai thực tế.
\end{itemize}


\subsubsection{Các tính năng liên quan đến đề tài}
Trong đồ án này, các khả năng sau của ESP32 được sử dụng trực tiếp:
\begin{itemize}
    \item \textbf{Giao tiếp I$^2$C:} Thực hiện kết nối với các cảm biến kỹ thuật số để thu thập dữ liệu nhiệt độ và độ ẩm.
    \item \textbf{GPIO/PWM:} Điều khiển trạng thái đèn LED, cấu hình hành vi của NeoPixel và các cơ cấu chấp hành khác.
    \item \textbf{Hỗ trợ máy chủ Web (Web server support):} Cho phép giám sát và điều khiển thiết bị cục bộ thông qua chế độ Access Point (AP).
    \item \textbf{Hỗ trợ hệ thống tệp trên Flash (ví dụ: LittleFS):} Lưu trữ các tài nguyên web tĩnh (HTML/CSS/JS).
    \item \textbf{Tương thích TinyML:} Khả năng suy luận trực tiếp trên thiết bị để phân tích trạng thái môi trường một cách nhẹ nhàng và tức thời.
\end{itemize}


\subsection{Cơ sở về RTOS}
FreeRTOS là một nhân hệ điều hành thời gian thực (RTOS) gọn nhẹ, mã nguồn mở, được thiết kế cho các hệ thống nhúng có tài nguyên hạn chế. Hệ điều hành này cung cấp một lớp trừu tượng có cấu trúc cho thực thi đồng thời, xử lý sự kiện và đồng bộ giữa các tác vụ mà không yêu cầu lập trình viên phải tự xây dựng cơ chế lập lịch phức tạp. Nhờ chi phí tài nguyên thấp và hệ sinh thái hỗ trợ trưởng thành, FreeRTOS được sử dụng rộng rãi trên các nền tảng vi điều khiển như ESP32, STM32 và ARM Cortex-M.

\subsubsection{Khái niệm cốt lõi}
\begin{itemize}
	\item \textbf{Task (Tác vụ):} là một đơn vị thực thi độc lập, thường được triển khai dưới dạng một vòng lặp vô hạn với không gian bộ nhớ (stack) và trạng thái thực thi riêng.
	\item \textbf{Scheduler (Bộ lập lịch):} thành phần cốt lõi quyết định tác vụ nào được thực thi tại mỗi thời điểm; FreeRTOS thường sử dụng cơ chế lập lịch ưu tiên trưng dụng (preemptive priority-based scheduling).
	\item \textbf{Queue (Hàng đợi):} cơ chế truyền thông an toàn luồng (thread-safe) để truyền tin nhắn hoặc dữ liệu giữa các tác vụ.
	\item \textbf{Semaphore:} một cơ chế đồng bộ hóa được sử dụng để báo hiệu các sự kiện và kiểm soát quyền truy cập vào các tài nguyên dùng chung.
	\item \textbf{Mutex:} một cơ chế khóa chuyên dụng để bảo vệ các tài nguyên dùng chung quan trọng khỏi việc bị truy cập đồng thời.
\end{itemize}


\subsubsection{Vai trò của FreeRTOS trong đồ án}
Trong hệ thống IoT này, FreeRTOS được dùng để điều phối nhiều mô-đun có yêu cầu khắt khe về mặt thời gian và truyền thông:
\begin{enumerate}
	\item \textbf{Quản lý tác vụ đồng thời:} các tác vụ riêng biệt đảm nhiệm việc xử lý web server, lấy mẫu cảm biến định kỳ, thực thi suy luận TinyML và điều khiển LED/thiết bị.
	\item \textbf{Đồng bộ hóa dữ liệu:} semaphore bảo vệ các luồng dữ liệu cảm biến dùng chung, trong khi hàng đợi hỗ trợ việc truyền dữ liệu an toàn giữa các tác vụ.
	\item \textbf{Phản hồi theo mức độ ưu tiên:} tác vụ xử lý yêu cầu web có thể được gán mức ưu tiên cao hơn để đảm bảo khả năng tương tác mượt mà với người dùng; các tác vụ thu thập cảm biến và suy luận được chạy ở mức ưu tiên trung bình một cách có kiểm soát.
	\item \textbf{Cô lập tài nguyên:} mỗi tác vụ có một vùng nhớ (stack) chuyên dụng riêng, giúp giảm thiểu hiện tượng tràn/nhiễu bộ nhớ và cải thiện độ ổn định của toàn hệ thống.
\end{enumerate}

\subsubsection{Kiến trúc tác vụ sử dụng trong đồ án}
Thiết kế phần mềm tuân theo mô hình tổ chức đa tác vụ của FreeRTOS:
\begin{itemize}
	\item \textbf{Task WebServer (mức ưu tiên cao):} xử lý các yêu cầu HTTP và các tương tác trên giao diện người dùng cục bộ.
	\item \textbf{Task Sensor Reader (mức ưu tiên trung bình):} thu thập dữ liệu đo nhiệt độ và độ ẩm định kỳ.
	\item \textbf{Task TinyML (mức ưu tiên trung bình):} thực thi quá trình suy luận trực tiếp trên thiết bị (on-device inference) để dự đoán trạng thái.
	\item \textbf{Task Điều khiển LED/Thiết bị:} cập nhật trạng thái đầu ra của LED/NeoPixel và các cơ cấu chấp hành liên quan.
	\item \textbf{Task Quản lý Wi-Fi:} quản lý kết nối mạng và logic tự động kết nối lại khi mất tín hiệu.
	\item \textbf{Task Nhàn rỗi - Idle Task (mức ưu tiên thấp nhất):} chạy khi không có tác vụ ứng dụng nào sẵn sàng thực thi.
\end{itemize}


\subsection{Cơ sở về TinyML}
TinyML (Tiny Machine Learning) tập trung vào việc triển khai suy luận học máy trên các thiết bị nhúng có tài nguyên hạn chế như vi điều khiển và cảm biến thông minh. Thay vì chuyển luồng dữ liệu cảm biến thô lên máy chủ đám mây để xử lý AI, TinyML thực hiện dự đoán trực tiếp ở biên. Cách tiếp cận này giúp giảm độ trễ, giảm chi phí năng lượng cho truyền thông, tăng cường riêng tư và cho phép hệ thống tiếp tục hoạt động ngay cả khi mạng không ổn định.

\subsubsection{Khái niệm cốt lõi}
\begin{itemize}
    \item \textbf{Suy luận (Inference):} Sử dụng mô hình đã được huấn luyện trước để dự đoán các mẫu dữ liệu đầu vào mới; TinyML chủ yếu tập trung vào quá trình suy luận thay vì huấn luyện trực tiếp trên thiết bị (on-device training).
    \item \textbf{Lượng tử hóa mô hình (Model quantization):} Giảm độ chính xác của các trọng số trong mô hình (ví dụ: từ float32 xuống int8 hoặc int16) nhằm giảm dung lượng bộ nhớ và cải thiện hiệu suất thực thi.
    \item \textbf{TensorFlow Lite / TensorFlow Lite Micro:} Các môi trường thực thi (runtimes) tối giản, được thiết kế để triển khai các mô hình nhỏ gọn trên phần cứng nhúng và thiết bị IoT.
    \item \textbf{Trích xuất đặc trưng (Feature extraction):} Chuyển đổi dữ liệu cảm biến thô thành các đặc trưng có ý nghĩa, phù hợp làm đầu vào cho mô hình.
    \item \textbf{Trí tuệ nhân tạo biên (Edge AI):} Thực hiện các tính toán AI ngay tại nguồn tạo ra dữ liệu (thiết bị biên) thay vì phụ thuộc vào cơ sở hạ tầng đám mây từ xa.
\end{itemize}

\subsubsection{Quy trình phát triển TinyML}
Quy trình triển khai thường bao gồm các bước sau:
\begin{enumerate}
    \item Huấn luyện mô hình ngoại tuyến bằng cách sử dụng các framework học máy trên Python và tập dữ liệu lịch sử.
    \item Chuyển đổi mô hình đã huấn luyện sang định dạng triển khai nhẹ (ví dụ: \texttt{.tflite} hoặc các định dạng biểu diễn nhỏ gọn tương đương).
    \item Áp dụng các kỹ thuật tối ưu hóa và lượng tử hóa để giảm mức sử dụng bộ nhớ và tăng tốc độ thực thi.
    \item Tích hợp các tham số của mô hình vào firmware (ví dụ: dưới dạng tệp tin header của ngôn ngữ C).
    \item Triển khai lên vi điều khiển ESP32-S3 và đánh giá độ trễ thực thi, mức tiêu thụ bộ nhớ cùng chất lượng dự đoán.
\end{enumerate}

\subsubsection{Ứng dụng TinyML trong đồ án}
Trong hệ thống IoT này, TinyML được sử dụng để phân loại trạng thái môi trường từ dữ liệu cảm biến thời gian thực:
\begin{itemize}
    \item \textbf{Đặc trưng đầu vào (Input features):} Nhiệt độ ($^\circ$C) và độ ẩm (\%) thu thập từ cảm biến DHT20.
    \item \textbf{Lớp đầu ra (Output classes):} Các trạng thái hoạt động của môi trường (ví dụ: trạng thái không hợp lệ, chế độ chờ, cần làm mát, cần hút ẩm, cần tạo ẩm).
    \item \textbf{Lựa chọn mô hình:} Thuật toán Naive Bayes với phân phối chuẩn (Gaussian Naive Bayes - GNB), được huấn luyện trên các cặp mẫu nhiệt độ - độ ẩm lịch sử đã được gán nhãn.
    \item \textbf{Tích hợp Firmware:} Các tham số của mô hình được xuất ra tệp \texttt{naive\_bayes\_model.h} và biên dịch cùng với ứng dụng trên ESP32-S3.
\end{itemize}

Để giảm thiểu băng thông giao tiếp, hệ thống có thể chỉ gửi đi các kết quả suy luận ngắn gọn (nhãn trạng thái) thay vì liên tục truyền toàn bộ dữ liệu thô.

\subsection{Cơ sở lý thuyết mô hình Naive Bayes và Phân phối chuẩn (GaussianNB)}

\subsubsection{Định lý Bayes và thuật toán Naive Bayes}
Naive Bayes là một thuật toán phân loại có giám sát dựa trên nền tảng của định lý xác suất Bayes. Mục tiêu của thuật toán là tính toán xác suất để một điểm dữ liệu với vector đặc trưng $X = (x_1, x_2, ..., x_n)$ thuộc về một lớp $y$ cụ thể. Phương trình cốt lõi của định lý Bayes được biểu diễn như sau:

\begin{equation}
P(y|X) = \frac{P(X|y)P(y)}{P(X)}
\end{equation}

Trong đó:
\begin{itemize}
    \item $P(y|X)$ (\textbf{Posterior - Xác suất hậu nghiệm}): Xác suất điểm dữ liệu thuộc lớp $y$ khi đã biết các đặc trưng $X$. Đây là giá trị mô hình cần tối đa hóa.
    \item $P(y)$ (\textbf{Prior - Xác suất tiên nghiệm}): Xác suất xuất hiện của lớp $y$ trong toàn bộ tập dữ liệu, được tính bằng tỷ lệ số mẫu của lớp $y$ trên tổng số mẫu.
    \item $P(X|y)$ (\textbf{Likelihood - Khả năng}): Xác suất quan sát được bộ đặc trưng $X$ nếu điểm dữ liệu thực sự thuộc lớp $y$.
    \item $P(X)$ (\textbf{Evidence - Chứng cứ}): Xác suất xuất hiện của bộ đặc trưng $X$. Do $P(X)$ là hằng số với mọi lớp $y$, ta có thể loại bỏ thành phần này trong quá trình so sánh để tối ưu khối lượng tính toán.
\end{itemize}

Điểm đặc trưng làm nên tên gọi "Naive" (Ngây thơ) của thuật toán là \textit{giả định độc lập có điều kiện} (conditional independence). Thuật toán giả định rằng sự xuất hiện của một đặc trưng $x_i$ hoàn toàn không liên quan đến sự xuất hiện của bất kỳ đặc trưng $x_j$ nào khác khi đã biết lớp $y$. Nhờ giả định này, thành phần Likelihood phức tạp được đơn giản hóa thành tích của các xác suất thành phần:

\begin{equation}
P(X|y) = \prod_{i=1}^{n} P(x_i|y)
\end{equation}

Kết hợp lại, nguyên tắc quyết định phân loại (decision rule) của Naive Bayes là chọn lớp $\hat{y}$ sao cho biểu thức sau đạt giá trị lớn nhất:

\begin{equation}
\hat{y} = \arg\max_{y} \left( P(y) \prod_{i=1}^{n} P(x_i|y) \right)
\label{eq:naive_bayes_base}
\end{equation}

\subsubsection{Biến đổi công thức với Gaussian Naive Bayes (GaussianNB)}
Trong bài toán giám sát môi trường, các đặc trưng đầu vào (Nhiệt độ, Độ ẩm) là các giá trị liên tục (continuous data) thay vì rời rạc (discrete data). Việc đếm tần suất xuất hiện của một con số nhiệt độ chính xác để tính $P(x_i|y)$ là bất khả thi. Để giải quyết vấn đề này, thuật toán \textbf{Gaussian Naive Bayes (GaussianNB)} được áp dụng.

GaussianNB giả định rằng các giá trị liên tục của một đặc trưng trong một lớp nhất định tuân theo \textbf{phân phối chuẩn} (Gaussian Distribution). Thay vì lưu trữ toàn bộ tập dữ liệu, mô hình chỉ cần ước lượng hai tham số thống kê cho mỗi đặc trưng $i$ trong mỗi lớp $y$:
\begin{itemize}
    \item $\mu_{y,i}$: Giá trị trung bình (Mean) của đặc trưng $i$ trong lớp $y$.
    \item $\sigma_{y,i}^2$: Phương sai (Variance) của đặc trưng $i$ trong lớp $y$.
\end{itemize}

Khi đó, thành phần xác suất $P(x_i|y)$ được thay thế bằng hàm mật độ xác suất (Probability Density Function - PDF) của phân phối chuẩn:

\begin{equation}
P(x_i|y) = \frac{1}{\sqrt{2\pi\sigma_{y,i}^2}} \exp\left(-\frac{(x_i - \mu_{y,i})^2}{2\sigma_{y,i}^2}\right)
\label{eq:gaussian_pdf}
\end{equation}

Thay phương trình (\ref{eq:gaussian_pdf}) vào nguyên tắc quyết định (\ref{eq:naive_bayes_base}), ta thu được công thức phân loại hoàn chỉnh của GaussianNB đối với dữ liệu liên tục:

\begin{equation}
\hat{y} = \arg\max_{y} \left( P(y) \prod_{i=1}^{n} \frac{1}{\sqrt{2\pi\sigma_{y,i}^2}} \exp\left(-\frac{(x_i - \mu_{y,i})^2}{2\sigma_{y,i}^2}\right) \right)
\label{eq:gaussian_nb_final}
\end{equation}

Phương trình (\ref{eq:gaussian_nb_final}) chính là nền tảng toán học để huấn luyện mô hình trên Python, từ đó trích xuất các ma trận $\mu$ và $\sigma^2$ nhằm triển khai thuật toán tính toán tối ưu trên vi điều khiển ESP32.

\subsection{Cơ sở về CoreIoT}

\subsubsection{Luồng làm việc trên đám mây và Đường truyền dữ liệu}
CoreIoT là nền tảng quản lý IoT thời gian thực. Hệ thống sử dụng MQTT - giao thức truyền thông Publish/Subscribe hạng nhẹ, đặc biệt tối ưu cho vi điều khiển ESP32. Luồng dữ liệu hai chiều được định tuyến như sau:

\begin{itemize}
    \item \textbf{Chiều lên (Telemetry):} Trạng thái môi trường và thông số ngưỡng được đóng gói thành định dạng JSON, xuất bản (Publish) định kỳ lên đám mây tại topic \texttt{v1/devices/me/telemetry}.
    \item \textbf{Chiều xuống (RPC Control):} Đám mây đóng vai trò Broker điều hướng lệnh điều khiển. ESP32 đăng ký lắng nghe (Subscribe) tại topic \texttt{v1/devices/me/rpc/request/+} để bắt sự kiện từ người dùng và phản hồi xuống phần cứng.
\end{itemize}

\subsubsection{Các thành phần giám sát}
Không gian điều khiển của hệ thống được tích hợp trên Dashboard thông qua hai nhóm thành phần chính:
\begin{itemize}
    \item \textbf{Thành phần hiển thị:} Gồm Thẻ số (Digital Cards) và Biểu đồ đường (Line Charts) để giám sát và lưu trữ chuỗi dữ liệu (time-series) của cảm biến.
    \item \textbf{Thành phần điều khiển:} Gồm thanh trượt (Slider) và công tắc (Switch) liên kết với các phương thức RPC (\texttt{setThre}, \texttt{setStateLED}). Điều này cho phép thay đổi logic điều khiển trực tiếp trên nền tảng Web.
\end{itemize}


\subsection{Tóm tắt chương}
Chương này đã trình bày quá trình thiết kế và triển khai phần mềm cho hệ thống IoT. Nhờ ứng dụng hệ điều hành FreeRTOS, ESP32 đã vận hành trên kiến trúc đa nhiệm gồm: module thu thập dữ liệu (DHT20, LCD), cảnh báo trực quan (LED, NeoPixel), Web Server thiết lập mạng và giao tiếp đám mây CoreIoT qua MQTT. 







